\section{Further research questions and a staged programme}
\label{sec:agenda}

The general target remains open for this implementation. The results above
suggest concrete next questions that can be tested or proved without
assuming the desired state bound. The order below reflects their dependence
on the present code and mathematics.

\subsection{Find the summands before they become large}
The Fitting pass already detects summands that graph connectivity misses.
Its cost currently competes with the time saved by sparser later matrices.
Can the scalar endomorphism space be maintained incrementally when a crossing
is added, instead of solved anew? A useful result would give a complete update
algorithm with cost polynomial in the changed blocks, together with a
criterion deciding when reusing an old splitting remains valid. Every reused
projection must still commute with the new differential.

The current bounded candidate search is deliberately incomplete. A candidate
endomorphism may be invertible or nilpotent on an indecomposable summand and
give no split, even if another endomorphism would work. A next implementation
could use primary decomposition of endomorphisms, or a proved decomposition
algorithm for finite-dimensional module representations. Its cost must be
compared with continuing the unsplit complex. The practical success criterion
is a reduction in complete measured work including decomposition overhead,
not the mere production of additional summands.

\subsection{Extend the class with a finite continuation observation}
The common-$\theta$ condition reduces the problem to the dual numbers. What
is the next useful class? Two independently active dot directions lead toward
\[
 \F[x,y]/(x^2,y^2),
\]
where the classification problem is substantially different. A carefully
scoped goal is to classify continuation observations of a specified family of
two-direction complexes at the rank-two threshold, rather than classify all
modules. The correct statement must quantify over every allowed suffix.
Agreement on sampled closures is evidence for a conjecture, not a valid
replacement rule.

Another useful target is a small tangle with two or more matching types and a
restricted set of saddle maps. The attached towers seen in actual scans
motivate this direction: whole-component common-$\theta$ recognition is too
restrictive on the measured corpus. A normal form for this larger class could
make the present quotient active on ordinary diagrams. Rational-tangle and
four-ended tangle constructions provide relevant prior structure
\cite{thompson,immersed}; their
hypotheses and grading conventions must be retained.

\subsection{Prove a structural bound on the unresolved intervals of a scan}
The checkpoint theorem reduces the target to explicit bounds on $w$ and $g$.
An immediate mathematical question is whether a natural family of validated
diagrams, extending the already tractable braid and rational-tangle families,
admits a constructible order with
\[
 g+2^{w/2}+w\log(w+1)=O((\log n)^2).
\]
A useful theorem must also bound the cost of finding that order. Merely
asserting the existence of a favorable representation would not complete the
algorithmic argument.

Corollary~\ref{cor:coefficient-alphabet} provides another concrete target:
logarithmic frontier width and $O(\log^2 n)$ checkpoint gaps suffice when
the checkpoint coefficients have dot degree at most two. Can a marked,
graded, or restricted-tangle structure force that degree bound after
normalization? The condition concerns actual checkpoint coefficients;
arbitrary intermediate algebra must still be charged in full.

A universal low-width replacement theorem for all knot types is unavailable:
Lunel--de Mesmay give torus-knot families with unavoidable large diagram
treewidth~\cite{widthobstruction}. Such lower bounds do not prevent an
algorithm from rejecting those knots using an inexpensive independent
certificate. This suggests an asymmetric target: either produce a short
nontriviality certificate, or construct a sequence of suitably bounded
algebraic checkpoints. Establishing this dichotomy would be far stronger
than optimizing a handful of observed scans.

\subsection{Build a catalogue of sharp obstructions}
The length-two bound has a sharp formal-block example. Similar sharpness
examples are needed for every proposed observation quotient. In particular,
one should construct common-nilpotent and mixed-matching families whose
serialized component sizes grow, and determine which are realizable inside
actual knot scans. A representation that is legal algebraically need not be
the Khovanov complex of a tangle. Keeping these two questions separate will
prevent synthetic compression demonstrations from being misreported as knot
benchmarks.

The 2026 explicit-basis exponential examples~\cite{kelomaki} are especially
valuable stress families. For each one, record how multiplicity, component
size, scalar endomorphism dimension, and checkpoint gaps grow after each new
normalization. The point is to identify a surviving obstruction, or prove
that an observation quotient removes it, rather than infer asymptotics from
a short timing curve.

\subsection{Exploit graded or marked structure with compatible gluing}
The present exact-rank backend retains homological degree but forgets quantum
degree; decision compression discards still more information. A separate
graded implementation could restrict scalar basis changes to compatible
bigraded blocks and use the grading to constrain possible endomorphisms.
This may lower the cost of the Fitting solver, but could also remove useful
ungraded decompositions. Both outcomes should be measured.

Cohen's characteristic-two splitting of arc algebras is relevant to reducing
the size of the coefficient algebra~\cite{cohen}. The tangle-bimodule result
does not automatically preserve both actions in arbitrary configurations.
A concrete next problem is to specify compatible marked endpoints throughout
the scan and prove that the chosen identifications commute with every gluing
operation. An algebra isomorphism alone does not supply that compatibility.

\subsection{Add a complete tractable topological class}
The linear-time recognition theorem for diagrams of treewidth two provides
an independent complete structural backend~\cite{treewidthtwo}. Implementing
its certified reductions would broaden the front end and offer a useful
comparison with the algebraic route. A proper implementation should accept a
verified decomposition, record every reduction, and distinguish recognition
on the promised class from attempting to discover a suitable diagram.

\subsection{Connect back to compressed hierarchy operations}
The geometric route remains valuable. If an implementation has at most
$L(g+1)^L$ phases and each phase costs at most $B$ bit operations, then
\[
 \log_2 T\le \log_2 B+\log_2 L+L\log_2(g+1)+O(1).
\]
Here the letters are the hierarchy parameters, not the checkpoint gap of
this article. This elementary accounting identity specifies what is needed
from the hierarchy programme: controlled length, controlled pattern data,
and bounded exact operations on the compressed encoding. The July 2026
preprint supplies important pieces~\cite{lackenby2026}; it does not justify
omitting any of these implementation costs.

\subsection{Formalize the new algebraic interface}
The continuation theorem is a manageable formalization target. It can be
split into finite-dimensional linear algebra over $\F$, finite interval
decomposition, the dual-number tensor pairing, and the parity refinement.
The largest foundational dependency is the derived classification; the
appendix provides an explicit elementary route that can be used instead of
importing a general derived-category classification theorem.

A practical first verification milestone is an independently checkable
certificate for each scalar basis split: binary change-of-basis matrices,
their inverses, and the transformed differential. For interval normalization,
rank certificates for the products in the multiplicity formula would certify
the inferred decomposition. Formalizing the observation-preserving shortening
rule then allows a small verifier to accept a decision model without trusting
the optimizer that found it.
